\SourceLeaf{030}per \(\Omega\), patet, \(2\xi,2\eta,2\zeta\) etc. esse quotientes differentiales partiales functionis \(\Omega\), puta
\[
2\xi=\frac{\mathrm{d}\Omega}{\mathrm{d}x},\quad2\eta=\frac{\mathrm{d}\Omega}{\mathrm{d}y},\quad2\zeta=\frac{\mathrm{d}\Omega}{\mathrm{d}z}\text{ etc.}
\]
Quapropter valores incognitarum ex observationum combinatione maxime idonea prodeuntes, quos \textit{valores maxime plausibiles} commode vocare possumus, identici erunt cum iis, per quos \(\Omega\) valorem minimum obtinet. Iam \(V-L\) indefinite exprimit differentiam inter valorem computatum et observatum. Valores itaque incognitarum maxime plausibiles iidem erunt, qui summam quadratorum differentiarum inter quantitatum \(V,V',V''\) etc. valores observatos et computatos, per observationum pondera multiplicatorum, minimam efficiunt, quod principium in \textit{Theoria Motus Corporum Coelestium} longe alia via stabiliveramus. Et si insuper praecisio relativa singularum determinationum assignanda est, per eliminationem indefinitam ex aequationibus (III) ipsas \(x,y,z\) etc. in tali forma exhibere oportet:
\[
\left.\begin{aligned}
x&=A+[\alpha\alpha]\xi+[\alpha\beta]\eta+[\alpha\gamma]\zeta+\text{etc.}\\
y&=B+[\beta\alpha]\xi+[\beta\beta]\eta+[\beta\gamma]\zeta+\text{etc.}\\
z&=C+[\gamma\alpha]\xi+[\gamma\beta]\eta+[\gamma\gamma]\zeta+\text{etc.}
\end{aligned}\right\}\quad\text{(VII)}
\]
\begin{center}etc.\end{center}
quo pacto valores maxime plausibiles incognitarum \(x,y,z\) etc. erunt resp. \(A,B,C\) etc., atque pondera his determinationibus tribuenda \(\frac1{[\alpha\alpha]},\frac1{[\beta\beta]},\frac1{[\gamma\gamma]}\) etc., sive errores medii in ipsis metuendi
\[
\begin{aligned}
\text{pro }x\mathord{.}\,\mathord{.}\,\mathord{.}\,\mathord{.}\,\mathord{.}\,\mathord{.}\,&m\BareRoot{p[\alpha\alpha]}=m'\BareRoot{p'[\alpha\alpha]}=m''\BareRoot{p''[\alpha\alpha]}\text{ etc.}\\
\text{pro }y\mathord{.}\,\mathord{.}\,\mathord{.}\,\mathord{.}\,\mathord{.}\,\mathord{.}\,&m\BareRoot{p[\beta\beta]}=m'\BareRoot{p'[\beta\beta]}=m''\BareRoot{p''[\beta\beta]}\text{ etc.}\\
\text{pro }z\mathord{.}\,\mathord{.}\,\mathord{.}\,\mathord{.}\,\mathord{.}\,\mathord{.}\,&m\BareRoot{p[\gamma\gamma]}=m'\BareRoot{p'[\gamma\gamma]}=m''\BareRoot{p''[\gamma\gamma]}\text{ etc.}
\end{aligned}
\]
\begin{center}etc.\end{center}
quod convenit cum iis, quae in \textit{Theoria Motus Corporum Coelestium} docuimus.

\Article{22}
De casu omnium simplicissimo, simul vero frequentissimo, ubi unica incognita adest, atque \(V=x,V'=x,V''=x\) etc., paucis seorsim agere conveniet. Erit scilicet \(a=\BareRoot{p},a'=\BareRoot{p'},a''=\BareRoot{p''}\) etc., \(l=-L\BareRoot{p},l'=-L'\BareRoot{p'},l''=-L''\BareRoot{p''}\) etc., et proin
\SourceLeaf{031}
\[
\xi=(p+p'+p''+\text{etc.})x-(pL+p'L'+p''L''+\text{etc.})
\]
Hinc porro
\[
\begin{aligned}
\relax[\alpha\alpha]&=\frac1{p+p'+p''+\text{etc.}}\\[6pt]
A&=\frac{pL+p'L'+p''L''+\text{etc.}}{p+p'+p''+\text{etc.}}
\end{aligned}
\]
Si itaque e pluribus observationibus inaequali praecisione gaudentibus, et quarum pondera resp. sunt \(p,p',p''\) etc., valor eiusdem quantitatis inventus est e prima \(=L\), e secunda \(=L'\), e tertia \(=L''\) etc., huius valor maxime plausibilis erit
\[
=\frac{pL+p'L'+p''L''+\text{etc.}}{p+p'+p''+\text{etc.}}
\]
pondusque huius determinationis \(=p+p'+p''\text{ etc.}\) Si omnes observationes aequali praecisione gaudent, valor maxime plausibilis erit
\[
=\frac{L+L'+L''+\text{etc.}}{\pi}
\]
i. e. aequalis medio arithmetico valorum observatorum, huiusque determinationis pondus \(=\pi\), accepto pondere observationum pro unitate.
\par\nobreak\medskip\centerline{\rule{.25\linewidth}{.3pt}}
